\subsection{Detector Simulation}
In this, so called, \emph{detsim} stage; after the effective production of ionisation electrons and propagation of scintillation photons, the signal formation is modelled. This stage takes as input the energy depositions (\texttt{edeps}) and the associated scintillation photons (\texttt{simphotons}) from the particle propagation step.  

\subsubsection{TPC Signal} 
\gls{wirecell} is a vast software toolkit that provides a comprehensive simulation of the TPC signal formation and reconstruction~\cite{wirecell}. It is defined as a tomographic automated 3D neutrino event reconstruction method for \glspl{lartpc}. Unlike paradigms that perform pattern recognition on 2D views first, \gls{wirecell} reconstructs the ionization density in 3D space prior to particle identification.

In \gls{dune}'s current simulation processes, the \texttt{edeps} are passed as a distributions of charge ($q_i$) at specific space-time points $(r_i, t_i)$ to produce a realistic detector responce. Conceptually, the simulation operates through five key functional stages that are intertwined in a convolutional chain~\cite{Adams_2018_I,Adams_2018_II}: 

\begin{enumerate}
    \item \textbf{Deposition (\texttt{edeps})}: Represents the initial distribution of ionization electrons in space and time.
    \item \textbf{Drift}: Models the transport of electrons to the anode, accounting for \textbf{attenuation} (due to impurities) and both \textbf{longitudinal and transverse diffusion}.
    \item \textbf{Duct (Field Response)}: This is a convolution of the position-dependent \textbf{field response function} ($F$) and the electronics response. It utilizes Garfield-simulated field responses based on Ramo's theorem to model current induction, including \textbf{long-range effects} on neighboring wires.
    \item \textbf{Digitization}: Models the front-end electronics based on specific sampling rates (e.g., 2 MHz), resolution, and dynamic range to produce ADC waveforms.
    \item \textbf{Noise}: Inherent electronics noise is added using a data-driven analytical model, often employing a Rayleigh-distributed amplitude spectrum.
\end{enumerate}

The \gls{wirecell} simulation has been rigorously validated against data from experiments such as \gls{microboone} and \gls{protodune-sp}. By modeling signal formation from first principles, it provides a crucial foundation for developing 3D reconstruction algorithms and evaluating systematic uncertainties in experiments like \gls{dune}.

To improve charge extraction, the waveforms undergo a \emph{deconvolution} procedure, where the electronics response and noise spectrum are inverted in frequency space to recover a cleaner signal representation. The deconvolved waveforms are evaluated under a \gls{roi} finder algorithm and trimmed for subsequent evaluation. This last step reduces the data volume by focusing only on regions with significant signal activity and trimming all information around the signal peak that is compatible with baseline noise.

\subsubsection{PDS Signal}
Simultaneously, the scintillation photons are detected by the \gls{pds}, composed of \gls{xa}s. The \texttt{simphotons} from Geant4 are converted into raw signals. Similarly to the \gls{tpc} reconstruction, the \gls{pds} response is also emulated with realistic timing resolution, dark noise, and electronic gain amplification. Simulated waveforms are produced for each optical channel and overlaid with noise. Also in this case, fixed \gls{roi} windows are defined after the waveforms surpass a certain threshold, indicating potential scintillation events. In contract to the \gls{tpc} \gls{roi}s, the waveforms of the \gls{pds} can be stored and process entirely due to the limited number of channels so entire \SI{16}{\micro\second} long \gls{roi}s are stored.

Once the number of arriving photons as a function of time is obtained per \gls{pd} channel, the digitization of the signal is performed. The generation of optical waveforms takes into account the expected \gls{pde} of the \gls{xa}s, the single \gls{pe} signal response, \gls{sipm} crosstalk, dark count rate, and electronic noise. Those waveforms go through a peak finding algorithm. For the \gls{vd} case, this is simply done by looking for peaks above a threshold (usually corresponding to 1.5 \gls{pe}s) and assigning the number of \gls{pe}s as the integral of the peak divided by the area corresponding to a single \gls{pe}. However, the \gls{hd} case presents additional complexity, as the single \gls{pe} shape shows a sizeable undershoot which requires a deconvolution to be performed before obtaining the number of \gls{pe}s in a given peak. 